% W1 outcome: a sharp G4 characterization/obstruction for (HM)_3.
% This is a fragment meant to use the notation and theorem environments of
% proof.tex.  It does not claim an unconditional proof of (HM)_3.

\subsection*{The third factorial moment: exact obstruction}

Write
\[
  \mathcal P_X=\{p\text{ prime}:X<p\le 2X\},\qquad M=X^2,
\]
where $X\ge 2$ is integral, and put
\[
  \Omega_p=\{0\le m<M:m\bmod p\in\mathcal Z_p\},\qquad
  f_p(m)=\mathbf 1_{\Omega_p}(m).
\]
Thus $K_X(m)=\sum_{p\in\mathcal P_X}f_p(m)$.  Set
\[
  A_p=|\Omega_p|,\qquad \rho_p=\frac{A_p}{M},\qquad
  \mu=\sum_{p\in\mathcal P_X}\rho_p
      =\frac1M\sum_{m<M}K_X(m),
\]
and retain
\[
  \lambda=\lambda_X=\sum_{p\in\mathcal P_X}\frac{Z(p)}p.
\]
For ordered distinct primes $p,q\in\mathcal P_X$, define
\begin{align*}
  J_{p,q}&=|\Omega_p\cap\Omega_q|,\\
  T_{p,q}&=\sum_{m\in\Omega_p\cap\Omega_q}
       \sum_{\ell\in\mathcal P_X\setminus\{p,q\}}f_\ell(m),\\
  \lambda_{p,q}&=\sum_{\ell\in\mathcal P_X\setminus\{p,q\}}
       \frac{Z(\ell)}\ell,\qquad
  \mu_{p,q}=\sum_{\ell\in\mathcal P_X\setminus\{p,q\}}\rho_\ell.
\end{align*}
The two aggregate positive pair--Palm excesses are
\begin{align}
  \mathscr P_3(X)
    &=\sum_{p\ne q}\bigl(T_{p,q}-J_{p,q}\lambda_{p,q}\bigr)_+,
       \label{eq:hm3-palm-periodic}\\
  \mathscr P_3^{\rm fin}(X)
    &=\sum_{p\ne q}\bigl(T_{p,q}-J_{p,q}\mu_{p,q}\bigr)_+.
       \label{eq:hm3-palm-finite}
\end{align}
The first uses the complete-period densities $Z(\ell)/\ell$; the second
uses the exact marginals on $[0,M)$ and therefore absorbs every one-column
boundary error.

\begin{theorem}[Sharp pair--Palm characterization of $(\mathrm{HM})_3$]
\label{thm:hm3-palm-characterization}
With all prime tuples ordered, one has the exact identities
\begin{align}
  \sum_{p\ne q}J_{p,q}
    &=\sum_{m<M}(K_X(m))_2,                              \label{eq:hm3-F2-id}\\
  \sum_{p\ne q}T_{p,q}
    &=\sum_{m<M}(K_X(m))_3.                              \label{eq:hm3-F3-id}
\end{align}
Moreover,
\begin{align}
  \mathscr P_3(X)
  &\le \sum_{m<M}(K_X(m))_3
  \le \mathscr P_3(X)+5M\lambda^3,                     \label{eq:hm3-palm-sandwich}\\
  \mathscr P_3^{\rm fin}(X)
  &\le \sum_{m<M}(K_X(m))_3
  \le \mathscr P_3^{\rm fin}(X)
       +5(1+2/X)M\lambda^3.                             \label{eq:hm3-palm-fin-sandwich}
\end{align}
Consequently $(\mathrm{HM})_3$ is equivalent to either one of
\[
  \mathscr P_3(X)\ll X^{2+o(1)}\lambda_X^3,
  \qquad
  \mathscr P_3^{\rm fin}(X)\ll X^{2+o(1)}\lambda_X^3.
\]
More generally, for every fixed $\theta\ge0$, the estimate in Goal~G2
with exponent $\theta$ is equivalent to either positive-excess estimate
with right-hand side $X^{2+\theta+o(1)}\lambda_X^3$.

There is also an exact centered formulation.  Define
\[
  V_X=\sum_{m<M}(K_X(m)-\mu)^2,
  \qquad
  C_{3,X}=\sum_{m<M}(K_X(m)-\mu)^3
\]
and the third factorial-moment defect
\begin{equation}
  \mathfrak D_3(X)
  =(C_{3,X}-M\mu)+3(\mu-1)(V_X-M\mu).                  \label{eq:hm3-defect-def}
\end{equation}
Then
\begin{equation}
  \sum_{m<M}(K_X(m))_3=M\mu^3+\mathfrak D_3(X).        \label{eq:hm3-defect-id}
\end{equation}
Thus $(\mathrm{HM})_3$ is equivalently the one-sided estimate
\[
  \bigl(\mathfrak D_3(X)\bigr)_+
  \ll X^{2+o(1)}\lambda_X^3.
\]
The same equivalence holds at every G2 exponent $\theta\ge0$.
\end{theorem}

\begin{proof}
If $p\ne q$ lie in $\mathcal P_X$, then $pq>M$.  Hence every residue
pair $(a\bmod p,b\bmod q)$ has at most one representative in $[0,M)$.
This removes, rather than estimates, the usual CRT boundary term.  At an
integer with $K_X(m)=k$, there are exactly $(k)_2$ ordered choices of
$(p,q)$, and each has $k-2$ choices of a third prime.  Its contribution
to $\sum_{p\ne q}T_{p,q}$ is therefore
$(k)_2(k-2)=(k)_3$.  This proves
\eqref{eq:hm3-F2-id}--\eqref{eq:hm3-F3-id}; in particular, no factor $3!$
is missing.

For nonnegative $A,B$,
\[
  (A-B)_+\le A\le B+(A-B)_+.
\]
Apply this with $A=T_{p,q}$ and $B=J_{p,q}\lambda_{p,q}$, sum over
ordered pairs, and use $\lambda_{p,q}\le\lambda$ together with the
proved second factorial moment:
\[
  \sum_{p\ne q}J_{p,q}\lambda_{p,q}
  \le\lambda\sum_{m<M}(K_X(m))_2
  \le5M\lambda^3.
\]
This proves~\eqref{eq:hm3-palm-sandwich}.

For each $p$ and each residue in $\mathcal Z_p$, the number of its
representatives below $M$ differs from $M/p$ by less than one.  Therefore
\[
  |A_p-MZ(p)/p|\le Z(p),
  \qquad
  |\mu-\lambda|\le\frac1M\sum_pZ(p)
       \le\frac{2\lambda}{X},                          \label{eq:hm3-mu-lambda}
\]
where the final inequality uses $p\le2X$.  In particular
$\mu\le(1+2/X)\lambda$.  Applying the same positive-part inequality
with $B=J_{p,q}\mu_{p,q}$ gives
\[
  \sum_{p\ne q}J_{p,q}\mu_{p,q}
  \le\mu\sum_{m<M}(K_X(m))_2
  \le5(1+2/X)M\lambda^3,
\]
which proves~\eqref{eq:hm3-palm-fin-sandwich}.  Each equivalence follows
in both directions from its sandwich.  This is a characterization, not a
new estimate hidden in a renamed hypothesis: the unproved arithmetic is
exactly the positive conditional overload of third primes on the actual
two-prime CRT loci.

Finally write $Y_m=K_X(m)-\mu$, so that $\sum_mY_m=0$.  Expanding
$K^3-3K^2+2K$ gives
\[
 \sum_{m<M}(K_X(m))_3
 =C_{3,X}+3(\mu-1)V_X+M(\mu^3-3\mu^2+2\mu),
\]
which rearranges to~\eqref{eq:hm3-defect-id}.  Since the main term is
nonnegative,
\[
  (\mathfrak D_3(X))_+
  \le\sum_{m<M}(K_X(m))_3
  \le M\mu^3+(\mathfrak D_3(X))_+.
\]
Equation~\eqref{eq:hm3-mu-lambda} bounds $M\mu^3$ by a constant multiple
of $M\lambda^3$, proving the last equivalences.  When $\lambda=0$ all
zero sets, all pair loci, and all displayed quantities vanish.  The same
is true of the Palm sums when there is no pair locus, so no division by a
possibly zero parameter has been suppressed.
\end{proof}

\subsection*{Why the available primewise and pairwise tools do not imply G2}

\begin{proposition}[An anchored rainbow star with bounded codegrees]
\label{prop:hm3-anchored-star}
For every sufficiently large integral $X$ there is a family
$(\mathcal Z_p^*)_{p\in\mathcal P_X}$ with the following properties.
Writing
\begin{align*}
 K_X^*(m)&=\sum_{p\in\mathcal P_X}
       \mathbf 1_{\{m\bmod p\in\mathcal Z_p^*\}},
 &\lambda_X^*&=\sum_{p\in\mathcal P_X}\frac{|\mathcal Z_p^*|}{p},\\
 \mu^*&=X^{-2}\sum_{m<X^2}K_X^*(m),
\end{align*}
one has:
\begin{enumerate}
\item Every $\mathcal Z_p^*$ has size zero or two.  Every nonempty set is
reflection-invariant, has no consecutive elements, avoids
$0,1,p-2,p-1$, and satisfies
\[
  |\mathcal Z_p^*\cap I|\le2|I|^{2/3}
\]
for every nonempty interval $I\subseteq\{0,\ldots,p-1\}$.
\item The universal CRT bound
\[
  \sum_{m<X^2}(K_X^*(m))_2\le5X^2(\lambda_X^*)^2
\]
holds, and even the global centered second moment has the natural scale
\[
  \sum_{m<X^2}(K_X^*(m)-\mu^*)^2\ll X^2\lambda_X^*.
\]
\item If $\mathcal P^*(m)=\{p:m\bmod p\in\mathcal Z_p^*\}$, then for all
distinct $0\le m,n<X^2$,
\[
  |\mathcal P^*(m)\cap\mathcal P^*(n)|\le7.
\]
Thus even a uniform $O(1)$ row-codegree bound does not prevent the model.
\item Nevertheless,
\[
  \lambda_X^*\asymp\frac1{\log X},
  \qquad
  \sum_{m<X^2}(K_X^*(m))_3
     \gg\frac{X^3}{(\log X)^3}.
\]
Consequently the family violates every G2 estimate with fixed
$\theta<1$.
\end{enumerate}
The family can additionally be extended to a reflection-invariant value
map whose largest fiber is at most two and whose collision energy is at
most $2p$.  Hence the numerical consequences of the bordered
nonvanishing theorem---the $O(p^{3/4})$ fiber bound and $O(p^{7/4})$
energy bound---also do not rule it out.

If one uses the gap-polynomial reflection identity
$N_h(-x-h-1)=(-1)^{h-1}N_h(x)$, the unique pair in every nonempty
$\mathcal Z_p^*$ also satisfies its corresponding necessary gap
certificate modulo $p$.
\end{proposition}

\begin{proof}
Put $M=X^2$, choose $m_0=\lfloor3M/4\rfloor$, and let $d=2m_0+1$.
Delete from $\mathcal P_X$ the primes dividing
\[
  E=(m_0-1)m_0(m_0+1)(m_0+2)d.
\]
Only $O(1)$ primes are deleted: $|E|\ll X^{10}$ and every deleted prime
exceeds $X$.  For every remaining prime define
\[
  a_p=m_0\bmod p,
  \qquad
  \mathcal Z_p^*=\{a_p,p-1-a_p\},
\]
and put $\mathcal Z_p^*=\varnothing$ for deleted primes.  The deletion
makes the two displayed residues distinct and excludes the four endpoint
residues.  Their difference is a nonzero even integer, so they are not
consecutive.  Reflection is built in, and the interval estimate is
immediate because an interval contains at most two marked residues.

Let $L$ be the number of nonempty columns.  The prime number theorem and
Mertens' theorem give
\[
  L\sim\frac{X}{\log X},
  \qquad
  \lambda_X^*=2\sum_{\substack{X<p\le2X\\p\nmid E}}\frac1p
  \asymp\frac1{\log X}.
\]
Every active prime marks $m_0$, so $K_X^*(m_0)=L$ and
\[
  \sum_{m<M}(K_X^*(m))_3\ge(L)_3
  \gg\frac{X^3}{(\log X)^3}.
\]
Dividing by $X^{2+\theta+o(1)}(\lambda_X^*)^3$ leaves
$X^{1-\theta-o(1)}$, which tends to infinity for $\theta<1$.
The second-factorial-moment estimate is universal and follows by the same
CRT argument as Proposition~\ref{prop:hm2}.  Also
\begin{align*}
  \sum_m(K_X^*(m)-\mu^*)^2
  &\le\sum_m(K_X^*(m))^2\\
  &=\sum_mK_X^*(m)+\sum_m(K_X^*(m))_2\\
  &\ll M\lambda_X^*+M(\lambda_X^*)^2
  \ll M\lambda_X^*.
\end{align*}

It remains to check the codegree.  Fix $0\le m<n<M$ and put $h=n-m$.
If an active prime marks both integers, their two residues are chosen from
$a_p$ and $-1-a_p$.  Equal choices give $p\mid h$; different choices give
\[
  h\equiv\pm(2a_p+1)\equiv\pm d\pmod p.
\]
Thus every common prime divides the nonzero integer
\[
  H=h(h-d)(h+d).
\]
Here $0<h<M<d$, and $|H|<8M^3=8X^6$.  If there are $k$ common primes,
their product exceeds $X^k$ and divides $|H|$, whence
\[
  k\le\frac{\log(8X^6)}{\log X}\le7
\]
for sufficiently large $X$.

For the value-map assertion, partition the residues into the orbits of
$r\mapsto p-1-r$.  At an active prime, assign zero to the marked orbit
and distinct nonzero values to all other orbits.  At a deleted prime,
assign distinct nonzero values to every orbit; this is possible because
there are $(p+1)/2\le p-1$ orbits.  Thus the zero fiber is exactly
$\mathcal Z_p^*$ in both cases.  Every fiber has size at most two, and,
since $t^2\le2t$ for $t\in\{0,1,2\}$, the fiber energy is at most $2p$.

For completeness, the gap-polynomial reflection identity used here is
symbolic.  For $h\ge2$, let $D_h(x)$ be the symmetric
$(h-1)\times(h-1)$ tridiagonal matrix with diagonal entries
$P(x+i)$, $1\le i<h$, and off-diagonal entries
$(x+i+1)^3$, $1\le i<h-1$.  Expansion in the last row and the defining
recurrence for $N_h$ give $\det D_h(x)=N_h(x)$.  If $R$ is the matrix
that reverses the coordinate order, then $P(-v-1)=-P(v)$ gives, entry by
entry,
\[
  D_h(-x-h-1)=-R D_h(x)R.
\]
Consequently
\begin{equation}
  N_h(-x-h-1)=(-1)^{h-1}N_h(x).                       \label{eq:hm3-gap-reflection}
\end{equation}
(For $h=1$ this is immediate from $N_1=1$.)  This proves the identity
for every $h$, without numerical verification.

Finally let $u=\min(a_p,p-1-a_p)$ and $h_p=p-1-2u$.
The deleted endpoints ensure $1\le u<u+h_p<p$; moreover $h_p$ is positive
and even.  Since
\[
  -u-h_p-1=u-p\equiv u\pmod p,
\]
the gap-polynomial reflection identity gives
$N_{h_p}(u)\equiv-N_{h_p}(u)\pmod p$, hence
$N_{h_p}(u)\equiv0\pmod p$ for odd $p$.
\end{proof}

\begin{remark}[Scope of the obstruction]
The preceding family is an incidence model, not the actual zero family of
the Ap\'ery recurrence.  It proves that reflection, no consecutive zeros,
the interval $O(|I|^{2/3})$ bound, the universal CRT second moment, the
primewise fiber/energy bounds, the gap-polynomial pair certificates, and
even uniform $O(1)$ row codegree do not imply any G2 power saving.  The
missing arithmetic must constrain, across distinct primes, which of the
structurally permitted reflection orbits is selected as the zero fiber.
The pair--Palm quantity in Theorem~\ref{thm:hm3-palm-characterization}
is the minimal aggregate positive overload that measures exactly this
selection problem.
\end{remark}

\subsection*{Computational findings}

The direct CRT enumerator in \texttt{hm3\_explore.py}, independently
cross-checked against an $m$-scatter, gives the following ordered third
factorial moments:
\[
\begin{array}{c|r|r|c}
X&\sum_{m<X^2}(K_X(m))_3&X^2\lambda_X^3&R_3\\ \hline
256&60&131.989926&0.454580146\\
512&150&209.308570&0.716645284\\
1024&486&583.330960&0.833146247\\
2048&1824&1953.964920&0.933486564\\
4096&7218&7608.237134&0.948708600
\end{array}
\]
At $X=4096$, the $1203$ canonical prime triples are supported on $1161$
integers; $1147$ have $K_X=3$ and $14$ have $K_X=4$.  The conditional
third-prime extension rate per unordered pair hit is $0.073422$, close to
$\lambda_X=0.0768283$.  Quotient repetitions, reflection-orbit pileups,
and residue correlations show no visible clustering.  Full histograms and
reproduction details are in \texttt{hm3\_exploration.md}.  These data are
evidence for $(\mathrm{HM})_3$, not a proof.

\subsection*{Stall report and infrastructure audit}

Goals G1 and G2 are not proved here, and no G3 hypothesis is claimed.
In particular, a third-order Palm condition should not be advertised as
``weaker than AP--BDH'': the latter is an $L^2$ condition and the two
statements are not ordered without extra hypotheses.  The completed
deliverable is G4, Theorem~\ref{thm:hm3-palm-characterization}, together
with the sharpness obstruction above.

Two issues in the supplied route are material to this conclusion.
First, Lucas gives only
\[
  m\bmod p\in\mathcal Z_p\quad\Longrightarrow\quad p\mid b_m.
\]
The converse can fail through the leading digit.  For example, $b_2=73$;
with $X=64$, $p=73$, and $m=146<X^2$, Lucas gives
$b_{146}\equiv b_2b_0\equiv0\pmod{73}$, while
$m\bmod73=0\notin\mathcal Z_{73}$.  Hence
$K_X(m)\le\omega_{(X,2X]}(b_m)$, but equality and the corresponding
factorial-moment identity asserted in the specification are false.

Second, the written proof of Lemma~\ref{lem:codegree} starts with the
unavailable assertion $p>N/2\ge h$ and then says that the two residues
``differ by exactly $h$''.  Quotient wrap invalidates that sentence even
under the displayed extra inequalities: for $p=37$ one has
$\mathcal Z_{37}=\{17,19\}$, and $m=93$, $n=128$, $N=70$ give
$37>N/2\ge n-m$ while the residues are $19$ and $17$.  The lemma may be
repairable by a separate wrap argument using reflection, but the proof as
written cannot be used as an established HM3 input.  The anchored-star
proposition shows that even a repaired, stronger $O(1)$ row-codegree
estimate would still not control a one-row rainbow pileup.

Finally, Proposition~\ref{prop:profile} in the second-moment section states
$\|\rho\|_2^2\le T_N^2/N$, whereas Cauchy--Schwarz gives the reverse
inequality and the proposition's own table reports a ratio near $1.25$.
No step above uses that proposition.

The trivial tuple estimate used for comparison should also be read as the
uniform inequality
\[
  \sum_{m<X^2}(K_X(m))_3
  \le\Bigl(\sum_{X<p\le2X}Z(p)\Bigr)^3
  \le8X^3\lambda_X^3.
\]
It need not be asymptotic to $(2X\lambda_X)^3$ when the zero mass is
concentrated on fewer than three primes.
